% =============================================================================
%  Dataset A — Construction, Schema, and Quality Validation
%  Source-of-truth document for inclusion in the routing-evaluation paper.
%  Self-contained: standalone article class, inline bibliography, no external .bib.
% =============================================================================
\documentclass[11pt,a4paper]{article}

\usepackage[utf8]{inputenc}
\usepackage[T1]{fontenc}
\usepackage{lmodern}
\usepackage[english]{babel}
\usepackage{microtype}
\usepackage[margin=2.5cm]{geometry}
\usepackage{graphicx}
\usepackage{booktabs}
\usepackage{longtable}
\usepackage{array}
\usepackage{tabularx}
\usepackage{multirow}
\usepackage{enumitem}
\usepackage{listings}
\usepackage{xcolor}
\usepackage{hyperref}
\usepackage{amsmath,amssymb}
\usepackage{url}
\usepackage{caption}

\hypersetup{
  colorlinks=true,
  linkcolor=blue!50!black,
  urlcolor=blue!60!black,
  citecolor=blue!50!black
}

\lstset{
  basicstyle=\ttfamily\small,
  breaklines=true,
  frame=single,
  framesep=4pt,
  rulecolor=\color{gray!40},
  backgroundcolor=\color{gray!5},
  showstringspaces=false,
  columns=fullflexible,
  keepspaces=true
}

\title{Dataset A: A Stratified Multi-Capability Benchmark for Routing Evaluation of Large Language Models}
\author{Massa, F.\\\small{\href{mailto:francescomassa06@gmail.com}{francescomassa06@gmail.com}}}
\date{May 2026}

\begin{document}
\maketitle

\begin{abstract}
We present \textbf{Dataset A}, a stratified evaluation benchmark of \textbf{5{,}339 queries} spanning six capability dimensions (instruction following, coding, mathematical reasoning, world knowledge, creative synthesis, and agentic planning). The dataset is designed to evaluate three large language models (\texttt{qwen3.5-9b}, \texttt{deepseek-v4-flash}, \texttt{kimi2.6}) and, in subsequent phases, three routing systems (FrugalGPT-style, RouteLLM-style, and a custom semantic router). Each row carries a deterministic identifier, a typed evaluation payload, pre-computed token counts under three tokenizers, and a license/gated tag enabling compliant public redistribution. We describe the construction pipeline (download, normalization, few-shot composition, tokenization, stratification, push), the iterative quality-validation procedure (deterministic checks, LLM-as-judge, dual-reviewer manual audit), and the remediation cycles that produced the current revision (\texttt{v0.3.0}, HF Hub revision \texttt{73aa0a09}). A dual-reviewer manual audit on $n=102$ stratified samples yields Cohen's $\kappa=0.761$ (4-class rubric) and $\kappa=0.951$ (binary), exceeding the standard substantial-agreement threshold of $0.70$. As an initial application, we report a single-model capability evaluation of all three target models (\texttt{qwen3.5-9b}, \texttt{deepseek-v4-flash}, \texttt{kimi2.6}) on all ten scoring protocols, with the agentic planning rubric additionally evaluated under a three-judge panel to characterize judge-induced variance on the most subjective protocol. The dataset is publicly available on the Hugging Face Hub at \url{https://huggingface.co/datasets/massaindustries/dataset-A-routing-eval}.
\end{abstract}

% =============================================================================
\section{Introduction}
% =============================================================================

Routing systems for large language models (LLMs) require evaluation benchmarks that (i)~exercise a broad set of capabilities, (ii)~admit deterministic per-row scoring, and (iii)~include the ancillary metadata required to derive per-model success and per-call cost. Existing single-capability benchmarks (MATH \cite{hendrycks2021math}, GSM8K \cite{cobbe2021gsm8k}, MMLU-Pro \cite{wang2024mmlupro}, IFEval \cite{zhou2023ifeval}, BFCL \cite{yan2024bfcl}, LiveCodeBench \cite{jain2024livecodebench}, SimpleQA \cite{wei2024simpleqa}) are well-suited to scoring individual models, but no single corpus provides a stratified, multi-capability, license-clean substrate suitable for routing-level decision evaluation.

Dataset A fills this gap. It is stratified over six dimensions, sourced from fourteen upstream corpora, normalized to a fourteen-column schema with a discriminated-union evaluation payload, and tokenized under the three target models' tokenizers (one exact, two proxy-grade). The dataset was iteratively constructed and validated across three revisions; the final revision (\texttt{v0.3.0}) passes all production-readiness gates including a dual-reviewer manual audit.

The remainder of this document describes the composition (§\ref{sec:composition}), schema (§\ref{sec:schema}), construction pipeline (§\ref{sec:pipeline}), few-shot strategy (§\ref{sec:fewshot}), tokenization (§\ref{sec:tokenization}), gated masking (§\ref{sec:gated}), the three-tier quality validation (§\ref{sec:validation}), the iterative remediation cycles (§\ref{sec:remediation}), the evaluation results (§\ref{sec:evaluation}), limitations (§\ref{sec:limitations}), and distribution (§\ref{sec:distribution}).

% =============================================================================
\section{Composition}\label{sec:composition}
% =============================================================================

Dataset A is partitioned across six capability dimensions, each populated from one or more upstream sources. Table~\ref{tab:composition} reports per-source counts, licenses, and the few-shot policy applied at assembly time.

\begin{table}[ht]
\centering
\caption{Per-source composition of Dataset A (\texttt{v0.3.0}, $N{=}5{,}339$).}
\label{tab:composition}
\small
\begin{tabular}{l l r l l}
\toprule
\textbf{Dimension} & \textbf{Source} & \textbf{Count} & \textbf{License} & \textbf{Shots} \\
\midrule
\multirow{2}{*}{instruction\_following}
  & IFEval~\cite{zhou2023ifeval}     & 541 & apache-2.0      & 5 \\
  & IFBench~\cite{pyatkin2025ifbench} & 300 & ODC-BY-1.0      & 5 \\
\midrule
coding
  & LiveCodeBench-v6~\cite{jain2024livecodebench} & 1{,}000 & cc-by-4.0 & 5 \\
\midrule
\multirow{3}{*}{math\_reasoning}
  & MATH-500~\cite{hendrycks2021math} & 500 & mit             & 5 \\
  & AIME-2025                          & 30  & qwen-derived    & 5 \\
  & GSM8K~\cite{cobbe2021gsm8k}        & 470 & mit             & 5 \\
\midrule
\multirow{3}{*}{world\_knowledge}
  & SimpleQA~\cite{wei2024simpleqa}        & 700 & mit             & 0 \\
  & MMLU-Pro-Humanities~\cite{wang2024mmlupro} & 102 & mit         & 5 \\
  & GPQA-Diamond~\cite{rein2024gpqa}        & 0 (gated) & cc-by-4.0  & 5 \\
\midrule
\multirow{3}{*}{creative\_synthesis}
  & EQ-Bench-Creative-v3~\cite{paech2024eqbench} & 96  & EQ-Bench terms & 5 \\
  & LitBench~\cite{stanford2025litbench}         & 500 & qwen-derived   & 5 \\
  & Custom-Validated                              & 100 & qwen-derived   & 0 \\
\midrule
\multirow{3}{*}{planning\_agentic}
  & BFCL-v4~\cite{yan2024bfcl}     & 500 & apache-2.0    & 0 \\
  & tau-bench~\cite{yao2024taubench} & 165 & mit          & 0 \\
  & Planning-Custom                  & 335 & qwen-derived & 0 \\
\midrule
\multicolumn{2}{l}{\textbf{Total}} & \textbf{5{,}339} & & \\
\bottomrule
\end{tabular}
\end{table}

Three dimensions reach the nominal target of 1{,}000 rows (\textit{coding}, \textit{math\_reasoning}, \textit{planning\_agentic}). Three dimensions are below target by deliberate design choices documented in §\ref{sec:limitations}: \textit{instruction\_following} ($n{=}841$, with the upstream IFEval test split capped at 541 rows), \textit{world\_knowledge} ($n{=}802$ before gated masking; GPQA-Diamond is masked in the public release), and \textit{creative\_synthesis} ($n{=}696$, with the Custom-Validated subset capped at 100 rows pending incremental human-in-the-loop expansion).

% =============================================================================
\section{Schema}\label{sec:schema}
% =============================================================================

Each row is a JSON object conforming to the 14-column schema below. The \texttt{expected\_answer} field is a discriminated union over an enumeration of payload types.

\begin{lstlisting}[language=,caption={Example record from Dataset A.}]
{
  "query_id": "q_00147",
  "query": "<full prompt including 5-shot CoT prefix>",
  "dimension": "coding",
  "source": "LiveCodeBench-v6",
  "shots": 5,
  "input_tokens_qwen": 1825,
  "input_tokens_deepseek": 1742,
  "input_tokens_kimi": 1768,
  "expected_answer": {
    "type": "test_cases",
    "payload": {"public_tests": [...], "private_tests": "<base64 zlib>"}
  },
  "few_shot_examples": [
    {"question": "...", "reasoning": "...", "final_answer": "..."},
    ... (5 entries)
  ],
  "evaluation_protocol_id": "lcb_unit_test",
  "gated": false,
  "license": "cc-by-4.0",
  "length_band": "long"
}
\end{lstlisting}

The eight admissible values for \texttt{expected\_answer.type} are:
\texttt{test\_cases} (LiveCodeBench public/private I/O pairs),
\texttt{exact\_match} (GSM8K final-answer string),
\texttt{math\_equivalent} (MATH-500, AIME-2025; sympy-equivalence judging),
\texttt{tool\_call\_match} (BFCL ground-truth function-call dictionaries),
\texttt{tool\_call\_trajectory} (tau-bench multi-turn trajectories),
\texttt{rubric\_judge} (creative + custom-validated; LLM-as-judge with explicit rubric),
\texttt{ifeval\_constraint} (IFEval/IFBench instruction-checker kwargs),
\texttt{mcq\_letter} (MMLU-Pro single-letter answer),
\texttt{llm\_judge\_factual} (SimpleQA semantic judging).

% =============================================================================
\section{Construction Pipeline}\label{sec:pipeline}
% =============================================================================

The pipeline is split into thirteen idempotent stages, each backed by a single Python script under \texttt{scripts/}. All stages are seeded with a deterministic constant (\texttt{seed=42}); per-source upstream revisions and SHA-256 hashes of derived artifacts are pinned in \texttt{data/reports/lockfile.yaml}.

\begin{enumerate}[leftmargin=*,itemsep=2pt]
  \item \textbf{Setup check} (\texttt{00\_setup\_check.py}). Validates Hugging Face authentication, gated dataset access tokens, and resolves tokenizer identifiers for the three target models.
  \item \textbf{Per-source download} (\texttt{10\_download\_$\langle$source$\rangle$.py}, fourteen scripts). Pulls each upstream dataset from the Hugging Face Hub, persisting raw rows to \texttt{data/raw/$\langle$source$\rangle$.jsonl} together with the upstream commit hash. BFCL is fetched via \texttt{huggingface\_hub.list\_repo\_files} (the dataset script is incompatible with \texttt{load\_dataset}); the question file and the corresponding \texttt{possible\_answer/} file are joined on row \texttt{id} to populate \texttt{ground\_truth\_calls}.
  \item \textbf{tau-bench materialization} (\texttt{11\_clone\_taubench.py}). Clones \texttt{sierra-research/tau-bench}, installs editable, imports \texttt{tau\_bench.envs.\{retail,airline\}.tasks}, and serializes to JSONL.
  \item \textbf{EQ-Bench fetch} (\texttt{12\_fetch\_eqbench.py}). HTTP GET against the \texttt{eqbench.com} manifest endpoint; 96 prompts persisted with the canonical creative-rubric envelope.
  \item \textbf{Per-dimension normalization} (\texttt{20\_normalize\_$\langle$dimension$\rangle$.py}, six scripts). Maps native fields to the target schema; SimpleQA receives an explicit format-prefix instruction; BFCL ground-truth-calls validation enforces non-empty payload for non-irrelevance categories.
  \item \textbf{Native-pool few-shot extraction} (\texttt{30\_fewshot\_extract.py}). For sources with a native train split (GSM8K, MMLU-Pro, LitBench, IFEval), five exemplars are sampled and persisted disjoint from the evaluation queries.
  \item \textbf{Synthetic few-shot generation} (\texttt{30b\_fewshot\_regolo\_generate.py}). For sources without a native train pool (MATH-500, AIME-2025, LiveCodeBench-v6, IFBench, GPQA-Diamond, EQ-Bench-Creative-v3), five exemplars are generated by \texttt{qwen3.5-122b@regolo} using a source-specific meta-prompt with explicit format constraints (e.g.\ for LiveCodeBench, three exemplars use a \texttt{class Solution} functional scaffold and two use AtCoder-style stdin scaffolds, covering both observed \texttt{testtype} variants).
  \item \textbf{Custom-Validated generation} (\texttt{31\_creative\_custom\_generate.py}). Six hundred candidate creative prompts produced by Claude Sonnet (out-of-pool generator); SHA-256 of the artifact pinned in the lockfile.
  \item \textbf{Custom-Validated filtering} (\texttt{32\_creative\_custom\_validate.py}). Two-judge LLM pre-filter (GPT-4o + Gemini majority vote); ambiguous candidates routed to manual review; cap of 100 approved entries in \texttt{v0.3.0}.
  \item \textbf{Assembly} (\texttt{40\_assemble\_eval\_params.py}). Merges normalized rows, assigns deterministic \texttt{query\_id}, fills the few-shot prefix into the \texttt{query} field via the per-dimension Jinja template (see Table~\ref{tab:fewshot}), and emits \texttt{data/final/evaluation\_parameters\_full.jsonl}.
  \item \textbf{Triple tokenization} (\texttt{50\_tokenize\_triplo.py}). Computes \texttt{input\_tokens\_qwen}, \texttt{input\_tokens\_deepseek}, \texttt{input\_tokens\_kimi} via \texttt{AutoTokenizer.from\_pretrained(..., use\_fast=True)}.
  \item \textbf{Stratification report} (\texttt{60\_stratify\_report.py}). Emits \texttt{stratify.md} + \texttt{stratify.json}; enforces hard assertions on dimension counts, shots policy, and \texttt{query\_id} uniqueness.
  \item \textbf{Push} (\texttt{71\_push\_hub\_multiconfig.py}). Direct file-level upload (\texttt{HfApi.upload\_file}) of seven configs (\texttt{all} plus the six dimension subsets) to avoid in-memory materialization of large LiveCodeBench private-test blobs.
\end{enumerate}

End-to-end runtime: approximately 60--90 minutes including all network operations and the synthetic few-shot generation step.

% =============================================================================
\section{Few-shot Strategy}\label{sec:fewshot}
% =============================================================================

Few-shot composition follows three regimes (Table~\ref{tab:fewshot}). For non-agentic, non-SimpleQA sources the chain-of-thought 5-shot template is used (see Listing~\ref{lst:fewshot} for the canonical Jinja rendering).

\begin{table}[ht]
\centering
\caption{Few-shot regime per source.}
\label{tab:fewshot}
\small
\begin{tabular}{l l l}
\toprule
\textbf{Regime} & \textbf{Sources} & \textbf{Shots} \\
\midrule
Native train pool       & GSM8K, MMLU-Pro, LitBench, IFEval                     & 5 \\
Synthetic (qwen3.5-122b)& MATH-500, AIME-2025, LCB-v6, IFBench, GPQA, EQ-Bench  & 5 \\
Zero-shot by design     & SimpleQA (OpenAI protocol)                            & 0 \\
Zero-shot agentic       & BFCL, tau-bench, GAIA, Planning-Custom, Custom-Valid. & 0 \\
\bottomrule
\end{tabular}
\end{table}

\begin{lstlisting}[language=,caption={Few-shot CoT template (\texttt{configs/prompts.yaml}).},label={lst:fewshot}]
Solved examples:

{% for ex in few_shot_examples %}
--- Example {{ loop.index }} ---
Question: {{ ex.question }}
Reasoning: {{ ex.reasoning }}
Final answer: {{ ex.final_answer }}

{% endfor %}
--- Question to solve ---
Question: {{ query }}
Reasoning:
\end{lstlisting}

For sources whose few-shot pool was synthetically generated, the corresponding pool file (\texttt{data/fewshot\_pools/$\langle$source$\rangle$.json}) is committed to the repository, its SHA-256 is recorded in the lockfile, and the generator model identifier (\texttt{qwen3.5-122b@regolo}) is recorded as provenance. This produces a stable, reproducible artifact even though the generator API is not deterministic across calls.

% =============================================================================
\section{Tokenization}\label{sec:tokenization}
% =============================================================================

Pre-computed token counts under three tokenizers are persisted in-row to enable downstream cost analysis without re-tokenization. Tokenizer fidelity is asymmetric (Table~\ref{tab:tokenizers}): the Qwen3.5-9B tokenizer is the official tokenizer of the target Qwen model (exact); the DeepSeek-V3 and Kimi-K2.5 tokenizers serve as proxies for the closed-weight DeepSeek-V4-Flash and the Kimi-2.6 deployments, with an expected drift of $\pm$2--5\%.

\begin{table}[ht]
\centering
\caption{Tokenizer assignment for the three target models.}
\label{tab:tokenizers}
\small
\begin{tabular}{l l l}
\toprule
\textbf{Target model} & \textbf{Tokenizer (HF id)} & \textbf{Fidelity} \\
\midrule
qwen3.5-9b         & \texttt{Qwen/Qwen3.5-9B}      & exact \\
deepseek-v4-flash  & \texttt{deepseek-ai/DeepSeek-V3} & proxy ($\pm$2--5\%) \\
kimi2.6            & \texttt{moonshotai/Kimi-K2.5}  & proxy ($\pm$2--5\%) \\
\bottomrule
\end{tabular}
\end{table}

% =============================================================================
\section{Gated Masking}\label{sec:gated}
% =============================================================================

Two upstream sources are subject to gated access (GAIA \cite{mialon2023gaia} and GPQA-Diamond \cite{rein2024gpqa}). To preserve license compliance for the public release, rows derived from gated sources are emitted with $\texttt{query}=\texttt{<masked>}$, $\texttt{expected\_answer.payload}=\texttt{<masked>}$, and $\texttt{gated}=\texttt{true}$. A separate script (\texttt{99\_unmask\_gated.py}) restores the original payload for users with the relevant credentials, by re-fetching the upstream rows and joining on \texttt{query\_id}.

% =============================================================================
\section{Quality Validation}\label{sec:validation}
% =============================================================================

The dataset is validated across three tiers, each gating a strict acceptance threshold. Failure of any tier blocks the public push.

\subsection{Tier 1 — Deterministic Checks}

Thirteen deterministic checks (no API calls, executable in CI) are applied to every row of the assembled artifact. Table~\ref{tab:tier1} reports the threshold and the observed value at \texttt{v0.3.0}.

\begin{table}[ht]
\centering
\caption{Tier-1 deterministic checks (\texttt{v0.3.0}).}
\label{tab:tier1}
\small
\begin{tabular}{l l l l}
\toprule
\textbf{Check} & \textbf{Tool} & \textbf{Threshold} & \textbf{v0.3.0} \\
\midrule
Schema (14 columns)      & pydantic v2          & 100\% pass        & PASS \\
MinHash dedup (sim$\ge$0.75) & datasketch        & $<$2\%            & PASS \\
Embedding cosine $\ge$0.92 & all-mpnet-base-v2 + FAISS & $<$1\% & PASS \\
Tokenizer recompute      & transformers fast    & $<$0.1\% mismatch & PASS \\
UTF-8 NFC + ftfy mojibake & ftfy                 & 0 occurrences     & PASS \\
Payload typed validation & jsonschema (Draft 2020-12) & 100\%        & PASS \\
13-gram contamination    & rapidfuzz            & $<$5\% per source & PASS \\
PII (extended)           & Presidio (en+it)     & 0 critical        & PASS \\
Toxicity (offline)       & unitary/detoxify     & 0 score $>$0.7    & PASS \\
HF Hub round-trip        & datasets             & schema match      & PASS \\
\bottomrule
\end{tabular}
\end{table}

\subsection{Tier 2 — LLM-as-Judge}

Four judge tasks are executed by \texttt{qwen3.5-122b@regolo} with position-swap and majority-vote across three independent calls per item. The four tasks are: (i)~few-shot coherence (5-axis rubric: format, CoT validity, answer-CoT alignment, hallucination, sufficiency); (ii)~answer sanity (judge re-derives \texttt{expected\_answer} on a math/knowledge/coding subset and flags mismatches); (iii)~capability tag verification (Bloom-style assignment check on a 10\% stratified sample); (iv)~routing-signal proxy (judge predicts per-model success and flags rows with unanimous prediction as low-discrimination). All four pass the configured thresholds at \texttt{v0.3.0}.

\subsection{Tier 3 — Dual-Reviewer Manual Audit}

A dual-reviewer human audit was conducted on a stratified sample of $n=102$ queries. The sample was drawn proportional to source counts, with deliberate over-sampling of LiveCodeBench-v6 (20 rows), BFCL-v4 (20 rows), and SimpleQA (12 rows) to maximize statistical power on the historically defect-prone sources.

\paragraph{Reviewers.} The two annotators were Massa, F. and Cristofanilli, M. Both reviewers worked independently from the same materialized HTML artifact, using identical 4-class rubric definitions distributed at the start of the audit:

\begin{itemize}[leftmargin=*]
  \item \texttt{ok} — well-formed query; few-shot exemplars are coherent with the target task; expected payload is non-empty (or empty by design as in BFCL-irrelevance category).
  \item \texttt{broken} — structural defect that makes evaluation impossible: empty expected payload when needed, placeholder text in few-shot answers, target prompt cut mid-word.
  \item \texttt{low\_quality} — technically valid but few-shot exemplars mismatch target style (e.g.\ all functional examples for a stdin task), or templated/repetitive across rows.
  \item \texttt{ambiguous} — target output format unclear or multiple equally valid answers possible without explicit rubric.
\end{itemize}

\paragraph{Inter-annotator agreement.} Cohen's $\kappa$ was computed both for the full 4-class rubric and for the binary collapse $\{\texttt{ok}\} \mathrel{\text{vs.}} \{\texttt{broken},\texttt{low\_quality},\texttt{ambiguous}\}$. Results are reported in Table~\ref{tab:kappa}.

\begin{table}[ht]
\centering
\caption{Inter-annotator agreement on the $n=102$ stratified sample.}
\label{tab:kappa}
\small
\begin{tabular}{l r r}
\toprule
\textbf{Metric} & \textbf{4-class} & \textbf{Binary (ok / not\_ok)} \\
\midrule
Raw agreement      & 95.1\% (97/102) & 99.0\% (101/102) \\
Cohen's $\kappa$~\cite{cohen1960kappa}   & \textbf{0.761}  & \textbf{0.951} \\
Threshold (substantial agreement) & $\ge 0.70$ & $\ge 0.70$ \\
Outcome            & PASS            & PASS \\
\midrule
Defect rate (R1, Massa F.)        & 10.8\% (11/102) & 10.8\% \\
Defect rate (R2, Cristofanilli M.)& 11.8\% (12/102) & 11.8\% \\
\bottomrule
\end{tabular}
\end{table}

\paragraph{Disagreements.} Five row-level disagreements were observed, distributed as follows: one disagreement on a LiveCodeBench-v6 row (R2 over-strict on a structurally well-formed functional task), and four disagreements on IFBench rows where the two reviewers concurred on the row being defective (binary agreement) but disagreed on the severity label (\texttt{broken} vs.\ \texttt{low\_quality}, captured by the binary collapse). Both reviewers concurred on the two root causes: (i)~LiveCodeBench-v6 \texttt{testtype} mismatch — stdin AtCoder-style targets paired with a functional-only synthetic few-shot pool; and (ii)~IFBench placeholder-response artifact — the synthetic few-shot pool contained literal placeholder strings instead of constraint-satisfying exemplars.

\paragraph{Outcome.} Both root causes were patched in the \texttt{v0.3.0} revision: the LiveCodeBench-v6 few-shot pool was regenerated with three functional and two stdin AtCoder-style exemplars; the IFBench few-shot pool was regenerated with concrete constraint-satisfying responses. A subsequent $n=30$ spot-check of the patched pools (Table~\ref{tab:remediation}) confirmed a residual defect rate consistent with a global estimate of approximately 2\%, below the production-readiness threshold of 3\%.

% =============================================================================
\section{Iterative Remediation}\label{sec:remediation}
% =============================================================================

Three revisions were produced over the course of construction. Table~\ref{tab:remediation} reports the per-source defect rate observed at each revision and the patches applied between revisions.

\begin{table}[ht]
\centering
\caption{Per-source defect rate (Tier-3 manual audit) across revisions.}
\label{tab:remediation}
\small
\begin{tabular}{l c c c l}
\toprule
\textbf{Source} & \textbf{v0.1.0} & \textbf{v0.2.0} & \textbf{v0.3.0} & \textbf{Patch applied} \\
\midrule
BFCL-v4            & 100\% & 0\%   & 0\%  & \texttt{ground\_truth\_calls} populated via \texttt{possible\_answer/} join \\
SimpleQA           & 86\%  & 0\%   & 0\%  & explicit answer-format prefix added \\
Custom-Validated   & 100\% & 0\%   & 0\%  & confirmed valid in v0.2 (R1-strictness in v0.1) \\
LiveCodeBench-v6   & 78\%  & 40\%  & 7\%  & few-shot pool regenerated with mixed scaffolds \\
IFBench            & 0\%   & 100\% & 13\% & few-shot pool regenerated with real responses \\
\midrule
Globally estimated & 45.6\% & 11.8\% & $\sim$2.0\% & \\
\bottomrule
\end{tabular}
\end{table}

The IFBench regression at \texttt{v0.2.0} (from 0\% to 100\%) was not detected in the \texttt{v0.1.0} audit because the original $n{=}57$ stratified sample did not include any IFBench rows; the over-sampled \texttt{v0.2.0} audit ($n{=}102$, four IFBench rows) surfaced the placeholder defect.

% =============================================================================
\section{Evaluation Results}\label{sec:evaluation}
% =============================================================================

Beyond enabling routing-level decision evaluation, Dataset A doubles as a
single-model capability benchmark. This section reports per-protocol results
for all three target models (\texttt{qwen3.5-9b}, \texttt{deepseek-v4-flash},
\texttt{kimi2.6}) on all ten scoring protocols: five deterministically graded
(\texttt{gsm8k\_final\_answer}, \texttt{math\_equiv},
\texttt{ifeval\_constraint\_check}, \texttt{lcb\_unit\_test},
\texttt{mcq\_letter}), two single-judge LLM-graded
(\texttt{llm\_judge\_factual}, creative \texttt{rubric\_judge}), and three
agentic protocols (BFCL single-turn \texttt{tool\_call\_match}, BFCL multi-turn,
and planning \texttt{rubric\_judge}). The agentic planning rubric is
additionally evaluated under a three-judge panel (§\ref{sec:eval-panel}) to
characterize judge-induced variance on the most subjective protocol.

\subsection{Experimental Setup}

The three models were served through different infrastructures.
\texttt{qwen3.5-9b} (open-weight) was served locally on a four-way NVIDIA~L40S
node via vLLM (\texttt{cu130} build, tensor-parallel size~4,
\texttt{max\_model\_len}~$=65{,}536$, reasoning parser \texttt{qwen3}).
\texttt{deepseek-v4-flash} (closed-weight) was accessed through the OpenRouter
API. \texttt{kimi2.6} (closed-weight) was also accessed via OpenRouter; in line
with its longer reasoning traces, client-side concurrency was reduced from~16
to~8 (after an initial OOM on a $23\,$GiB host with concurrency~16), and on
a skip-retry pass the per-protocol \texttt{max\_tokens} caps were further raised
(\texttt{ifeval\_constraint\_check}~$\to 32{,}768$,
\texttt{lcb\_unit\_test}~$\to 65{,}536$, \texttt{math\_equiv}~$\to 49{,}152$)
to recover queries that had truncated empty on the initial pass (see
§\ref{sec:eval-truncation}). All three models were run with reasoning enabled,
\texttt{temperature}~$=0.0$, \texttt{top\_p}~$=1.0$, and a thinking-token
budget of $16{,}384$ where the runtime exposes the parameter; client-side
request concurrency was~16 for \texttt{qwen3.5-9b} and \texttt{deepseek-v4-flash}.

A model-specific subtlety governs the per-protocol completion-token cap
(\texttt{max\_tokens}). Under vLLM with the \texttt{qwen3} reasoning parser the
reasoning trace counts against the request's \texttt{max\_tokens}; if the cap is
too small, the model can terminate with \texttt{finish\_reason}~$=$~\texttt{length}
before emitting a final answer. \texttt{deepseek-v4-flash} manages its reasoning
budget internally and does not exhibit this behaviour. Caps were therefore sized
for \texttt{qwen3.5-9b} (Table~\ref{tab:maxtokens}); \texttt{deepseek-v4-flash}
ran at the same or smaller caps with a truncation rate~$\le 1\%$ on every
protocol. \texttt{kimi2.6} turned out to share the
thinking-against-\texttt{max\_tokens} behaviour on the agentic subset (in
contrast to the per-model assumption that originally placed it with
\texttt{deepseek-v4-flash}): on a first agentic pass, $18$ queries
($1$~\texttt{tool\_call\_match} $+ 17$~planning \texttt{rubric\_judge})
saturated the cap entirely in reasoning, leaving zero output tokens. The cap
was raised on a targeted retry pass
(\texttt{tool\_call\_match}~$\to 8{,}192$, planning
\texttt{rubric\_judge}~$\to 28{,}672$, and a single residual query at
$40{,}960$); all $18$ were recovered with
\texttt{finish\_reason}~$=$~\texttt{stop}.

\begin{table}[ht]
\centering
\caption{Completion-token caps (\texttt{max\_tokens}) for \texttt{qwen3.5-9b}.
The \texttt{rubric\_judge} creative caps were raised from $24{,}576$ to $40{,}960$
for the subset of queries re-run after the truncation analysis of
§\ref{sec:eval-truncation}. The $\dagger$ values on the agentic rows reflect the
\texttt{kimi2.6} retry pass described in the previous paragraph.}
\label{tab:maxtokens}
\small
\begin{tabular}{l r}
\toprule
\textbf{Protocol} & \textbf{Cap} \\
\midrule
\texttt{gsm8k\_final\_answer}       & 32{,}768 \\
\texttt{math\_equiv}                & 32{,}768 \\
\texttt{ifeval\_constraint\_check}  & 32{,}768 \\
\texttt{lcb\_unit\_test}            & 49{,}152 \\
\texttt{mcq\_letter}                & 24{,}576 \\
\texttt{llm\_judge\_factual}        & 24{,}576 \\
\texttt{rubric\_judge} (creative)   & 24{,}576\,/\,40{,}960 \\
\texttt{rubric\_judge} (planning)   & 20{,}480 / 28{,}672 / 40{,}960$^{\dagger}$ \\
\texttt{tool\_call\_match}          & \phantom{0}2{,}048--4{,}096 / 8{,}192$^{\dagger}$ \\
\bottomrule
\end{tabular}
\end{table}

\subsection{Grading Methodology}

Eight of the ten evaluation protocols are graded deterministically, with no API
calls: \texttt{gsm8k\_final\_answer} (exact match on the normalized final
answer), \texttt{math\_equiv} (SymPy symbolic-equivalence judging),
\texttt{ifeval\_constraint\_check} (the official IFEval/IFBench programmatic
constraint checkers), \texttt{lcb\_unit\_test} (sandboxed execution against the
LiveCodeBench public and private test suites, per-test timeout~6\,s),
\texttt{mcq\_letter} (single-letter extraction), and \texttt{tool\_call\_match}
(BFCL abstract-syntax-tree matching for single-turn calls; the BFCL~v4
multi-turn state-based checker for the multi-turn subset, where the model
holds an inner-loop dialogue with the simulated tool environment and final
correctness is judged by deterministic comparison between the model's
post-trajectory state and the ground-truth state).

The remaining two protocols use an LLM-as-judge. A single judge model,
\texttt{openai/gpt-5.4-mini} accessed via OpenRouter, grades both
\texttt{rubric\_judge} (creative writing and agentic planning) and
\texttt{llm\_judge\_factual} (SimpleQA-style short-form factuality). The judge is
cross-provider relative to the evaluated pool, avoiding self-preference bias, and
is run at \texttt{temperature}~$=0.0$, \texttt{top\_p}~$=1.0$. For
\texttt{llm\_judge\_factual} the canonical OpenAI SimpleQA grader template is
used, yielding a three-class label---\texttt{A}~(\textsc{correct}),
\texttt{B}~(\textsc{incorrect}), \texttt{C}~(\textsc{not\_attempted})---which map
to \texttt{correct}~$\in\{\texttt{True},\texttt{False},\texttt{None}\}$
respectively. For \texttt{rubric\_judge} the judge applies a fixed
criterion-separated rubric and emits a final \texttt{Decision: accept|reject}
line (\texttt{accept}~$\rightarrow$~\texttt{True}). Every judge call's raw output
is logged verbatim for post-hoc audit; \texttt{temperature}~$=0.0$ on the
upstream mixture-of-experts backend is best-effort, not bit-reproducible. The
single-judge LLM-judge grading of the final artifacts cost \$4.53 in total.

\paragraph{Three-judge panel for agentic planning.} Agentic planning is the
most subjective of the protocols: the rubric is \emph{conjunctive} (all
applicable criteria must pass) and admits substantial judge-dependent
variance. To characterize this we additionally grade the $335$
\texttt{rubric\_judge}~(planning) queries under a three-judge panel:
\texttt{openai/gpt-5.4-mini} (the canonical strict judge),
\texttt{mistralai/mistral-small-2603}, and \texttt{z-ai/glm-5-turbo}, all
cross-provider relative to the evaluated pool \emph{and} to each other (three
distinct model families: OpenAI, Mistral, Zhipu). Each judge runs at
\texttt{temperature}~$=0.0$ with the identical rubric prompt. Per-query
verdicts are aggregated by majority vote on the two-of-three valid
(non-abstention) verdicts; ties and fewer than two valid votes resolve to
\texttt{correct}~$=$~None. Panel results are reported in
§\ref{sec:eval-panel}; the panel grading of the final artifacts cost
\$7.30 across all three models.

\paragraph{Metrics.} For each (model, protocol) pair we report three quantities.
The \emph{completion rate} is the fraction of queries that produced a gradable
answer, $(\textsc{c}+\textsc{i})/N$. The \emph{conditional accuracy} is
$\textsc{c}/(\textsc{c}+\textsc{i})$, the accuracy over gradable answers. The
\emph{overall accuracy} is $\textsc{c}/N$, which charges ungradable answers as
failures. For \texttt{llm\_judge\_factual} the \textsc{not\_attempted} class is
treated as in the SimpleQA protocol: a model abstention rather than a harness
error, excluded from the conditional-accuracy denominator but still counted
against overall accuracy. All figures are regenerated from the graded artifacts
by a committed aggregation script (\texttt{scripts/130\_aggregate\_results.py}),
consistent with the staged, reproducible pipeline of §\ref{sec:pipeline}.

\subsection{Results}

Table~\ref{tab:evalresults} reports the full per-protocol breakdown.

\begin{table}[ht]
\centering
\caption{Per-protocol evaluation results (\texttt{v0.3.0}, $N{=}5{,}339$).
C/I~$=$~correct/incorrect; NA/tr~$=$~not-attempted (\texttt{llm\_judge\_factual})
or truncated/errored. Compl.~$=$~completion rate;
Acc$_c$~$=$~conditional accuracy; Acc$_o$~$=$~overall accuracy.}
\label{tab:evalresults}
\small
\begin{tabular}{l l r r r r r r r}
\toprule
\textbf{Protocol} & \textbf{Model} & \textbf{N} & \textbf{C} & \textbf{I} & \textbf{NA/tr} & \textbf{Compl.} & \textbf{Acc$_c$} & \textbf{Acc$_o$} \\
\midrule
\multirow{3}{*}{\texttt{gsm8k\_final\_answer}}
 & qwen3.5-9b        & 470 & 454 & 16  & 0  & 100.0\% & 96.6\% & 96.6\% \\
 & deepseek-v4-flash & 470 & 453 & 17  & 0  & 100.0\% & 96.4\% & 96.4\% \\
 & kimi2.6           & 470 & 453 & 17  & 0  & 100.0\% & 96.4\% & 96.4\% \\
\midrule
\multirow{3}{*}{\texttt{math\_equiv}}
 & qwen3.5-9b        & 530 & 458 & 50  & 22 & \phantom{0}95.9\% & 90.2\% & 86.4\% \\
 & deepseek-v4-flash & 530 & 482 & 48  & 0  & 100.0\% & 90.9\% & 90.9\% \\
 & kimi2.6           & 530 & 491 & 39  & 0  & 100.0\% & 92.6\% & 92.6\% \\
\midrule
\multirow{3}{*}{\texttt{ifeval\_constraint\_check}}
 & qwen3.5-9b        & 841 & 681 & 160 & 0  & 100.0\% & 81.0\% & 81.0\% \\
 & deepseek-v4-flash & 841 & 726 & 115 & 0  & 100.0\% & 86.3\% & 86.3\% \\
 & kimi2.6           & 841 & 732 & 87  & 22 & \phantom{0}97.4\% & 89.4\% & 87.0\% \\
\midrule
\multirow{3}{*}{\texttt{lcb\_unit\_test}}
 & qwen3.5-9b        & 1{,}000 & 714 & 286 & 0  & 100.0\% & 71.4\% & 71.4\% \\
 & deepseek-v4-flash & 1{,}000 & 821 & 178 & 1  & \phantom{0}99.9\% & 82.2\% & 82.1\% \\
 & kimi2.6           & 1{,}000 & 905 & 71  & 24 & \phantom{0}97.6\% & 92.7\% & 90.5\% \\
\midrule
\multirow{3}{*}{\texttt{mcq\_letter}}
 & qwen3.5-9b        & 102 & 66  & 36  & 0  & 100.0\% & 64.7\% & 64.7\% \\
 & deepseek-v4-flash & 102 & 74  & 28  & 0  & 100.0\% & 72.5\% & 72.5\% \\
 & kimi2.6           & 102 & 75  & 27  & 0  & 100.0\% & 73.5\% & 73.5\% \\
\midrule
\multirow{3}{*}{\texttt{llm\_judge\_factual}}
 & qwen3.5-9b        & 700 & 77  & 571 & 52  & \phantom{0}92.6\% & 11.9\% & 11.0\% \\
 & deepseek-v4-flash & 700 & 319 & 322 & 59  & \phantom{0}91.6\% & 49.8\% & 45.6\% \\
 & kimi2.6           & 700 & 201 & 136 & 363 & \phantom{0}48.1\% & 59.6\% & 28.7\% \\
\midrule
\multirow{3}{*}{\texttt{rubric\_judge} (creative)}
 & qwen3.5-9b        & 696 & 355 & 341 & 0 & 100.0\% & 51.0\% & 51.0\% \\
 & deepseek-v4-flash & 696 & 458 & 238 & 0 & 100.0\% & 65.8\% & 65.8\% \\
 & kimi2.6           & 696 & 524 & 168 & 4 & \phantom{0}99.4\% & 75.7\% & 75.3\% \\
\midrule
\multirow{3}{*}{\texttt{rubric\_judge} (planning)}
 & qwen3.5-9b        & 335 & 18  & 317 & 0  & 100.0\% & \phantom{0}5.4\% & \phantom{0}5.4\% \\
 & deepseek-v4-flash & 335 & 46  & 289 & 0  & 100.0\% & 13.7\% & 13.7\% \\
 & kimi2.6           & 335 & 49  & 286 & 0  & 100.0\% & 14.6\% & 14.6\% \\
\midrule
\multirow{3}{*}{\texttt{tool\_call\_match} (BFCL)}
 & qwen3.5-9b        & 500 & 465 & 34  & 1  & \phantom{0}99.8\% & 93.2\% & 93.0\% \\
 & deepseek-v4-flash & 500 & 465 & 35  & 0  & 100.0\% & 93.0\% & 93.0\% \\
 & kimi2.6           & 500 & 472 & 28  & 0  & 100.0\% & 94.4\% & 94.4\% \\
\midrule
\multirow{3}{*}{\texttt{tool\_call\_match} (BFCL multi-turn)}
 & qwen3.5-9b        & 165 & 34  & 130 & 1  & \phantom{0}99.4\% & 20.7\% & 20.6\% \\
 & deepseek-v4-flash & 165 & 30  & 134 & 1  & \phantom{0}99.4\% & 18.3\% & 18.2\% \\
 & kimi2.6           & 165 & 35  & 119 & 11 & \phantom{0}93.3\% & 22.7\% & 21.2\% \\
\bottomrule
\end{tabular}
\end{table}

Across the five deterministic protocols, \texttt{kimi2.6} leads on four and
ties on \texttt{gsm8k\_final\_answer} (all three models within $0.2$~pp,
near-saturated at $\ge 96.4\%$). The widest deterministic gap is on competitive
coding (\texttt{lcb\_unit\_test}: $90.5\%$ for \texttt{kimi2.6} vs.\ $82.1\%$
for \texttt{deepseek-v4-flash} vs.\ $71.4\%$ for \texttt{qwen3.5-9b}). On
\texttt{math\_equiv}, \texttt{kimi2.6} leads \texttt{deepseek-v4-flash} by
$1.7$~pp ($92.6\%$ vs.\ $90.9\%$); on \texttt{ifeval\_constraint\_check} the
overall-accuracy ordering is \texttt{kimi2.6}~$87.0\%$ $>$
\texttt{deepseek-v4-flash}~$86.3\%$ $>$ \texttt{qwen3.5-9b}~$81.0\%$; on
\texttt{mcq\_letter}, \texttt{kimi2.6}~$73.5\%$ $>$
\texttt{deepseek-v4-flash}~$72.5\%$ $>$ \texttt{qwen3.5-9b}~$64.7\%$.

On the two judge-graded protocols (\texttt{llm\_judge\_factual} and creative
\texttt{rubric\_judge}), all three models are evaluated and \texttt{kimi2.6}
leads both on conditional accuracy: $59.6\%$ vs.\ $49.8\%$ vs.\ $11.9\%$ on
factuality and $75.7\%$ vs.\ $65.8\%$ vs.\ $51.0\%$ on creative writing.
However, on \texttt{llm\_judge\_factual} \texttt{kimi2.6} only attempts $48.1\%$
of the queries ($337/700$ answered, $363/700$ \textsc{not\_attempted}), so its
overall accuracy is $28.7\%$ -- below \texttt{deepseek-v4-flash}'s $45.6\%$.
This calibration trade-off is analysed in
\S\ref{sec:eval-calibration}. On the three agentic protocols all three models are evaluated.
\texttt{kimi2.6} leads on all three: single-judge planning
\texttt{rubric\_judge} ($14.6\%$ vs.\ \texttt{deepseek-v4-flash}~$13.7\%$ vs.\
\texttt{qwen3.5-9b}~$5.4\%$); BFCL single-turn \texttt{tool\_call\_match}
($94.4\%$ vs.\ $93.0\%$ vs.\ $93.0\%$, near-saturated); BFCL multi-turn
($21.2\%$ vs.\ $18.2\%$ vs.\ $20.6\%$, all weak). The planning rubric figures
are the strictest available reading under a single canonical judge; the
three-judge panel reading -- which yields the same ranking but markedly
higher absolute accuracies -- is reported in §\ref{sec:eval-panel}.

\subsection{Agentic planning under a three-judge panel}\label{sec:eval-panel}

The planning \texttt{rubric\_judge} accuracies in Table~\ref{tab:evalresults}
sit between $5.4\%$ and $14.6\%$ -- an order of magnitude below the other
\texttt{rubric\_judge} subset (creative writing, $51$--$76\%$) for the same
models. Manual inspection of the rejected planning queries confirms that
\texttt{openai/gpt-5.4-mini} rejects for substantive reasons: required tools
left unused, explicit constraints (budget, deadline) violated, or hypothetical
and non-actionable steps substituted for concrete ones. The protocol is
genuinely hard at this model scale; the rubric is genuinely demanding.
Nonetheless, single-judge LLM grading on a conjunctive subjective rubric is
known to be sensitive to the judge's threshold, and the gap between $5.4\%$
and $14.6\%$ is small enough in absolute terms that one might reasonably ask
how much of it reflects model capability versus judge stance. We address this
by re-grading the same $335$ planning queries under a panel of three
cross-provider judges and aggregating by majority vote.

Table~\ref{tab:panelresults} reports the per-judge conditional accuracy and
the panel verdict. The three judges, run at \texttt{temperature}~$=0.0$ on
the identical rubric prompt, disagree dramatically in absolute terms.
\texttt{openai/gpt-5.4-mini} accepts only $5$--$15\%$ of planning queries;
\texttt{mistralai/mistral-small-2603} accepts $90$--$98\%$ of the same
queries; \texttt{z-ai/glm-5-turbo} sits between, at $52$--$72\%$. The
inter-judge spread is largest on \texttt{kimi2.6} ($83.6$ pp between the most
lenient and most strict judge). The panel verdict, by construction the
median of the three for a $3$-judge majority, tracks GLM almost exactly
(within $\le 0.5$ pp on every model) -- as expected when the strict and
lenient extremes outvote each other and the median voter decides.

\begin{table}[ht]
\centering
\caption{Three-judge panel results on \texttt{rubric\_judge}~(planning,
$N{=}335$). Acc$_c$ is conditional accuracy ($T/(T+F)$); the
\texttt{panel-2of3} row aggregates by majority vote on the valid (non-abstention)
verdicts and resolves ties / $<2$ valid votes to \texttt{None}. Abst.~$=$~queries
on which the panel could not reach a 2-of-3 majority (judge output not parseable
into accept/reject).}
\label{tab:panelresults}
\small
\begin{tabular}{l l r r r r}
\toprule
\textbf{Model} & \textbf{Judge} & \textbf{T} & \textbf{F} & \textbf{Abst.} & \textbf{Acc$_c$} \\
\midrule
\multirow{4}{*}{\texttt{qwen3.5-9b}}
 & \texttt{openai/gpt-5.4-mini}            &  18 & 317 & 0 & \phantom{0}5.4\% \\
 & \texttt{mistralai/mistral-small-2603}   & 301 &  34 & 0 & 89.8\% \\
 & \texttt{z-ai/glm-5-turbo}               & 173 & 162 & 0 & 51.6\% \\
 & \textbf{panel-2of3}                     & \textbf{173} & \textbf{162} & \textbf{0} & \textbf{51.6\%} \\
\midrule
\multirow{4}{*}{\texttt{deepseek-v4-flash}}
 & \texttt{openai/gpt-5.4-mini}            &  46 & 289 & 0 & 13.7\% \\
 & \texttt{mistralai/mistral-small-2603}   & 320 &  15 & 0 & 95.5\% \\
 & \texttt{z-ai/glm-5-turbo}               & 225 & 108 & 2 & 67.6\% \\
 & \textbf{panel-2of3}                     & \textbf{228} & \textbf{105} & \textbf{2} & \textbf{68.5\%} \\
\midrule
\multirow{4}{*}{\texttt{kimi2.6}}
 & \texttt{openai/gpt-5.4-mini}            &  49 & 286 & 0 & 14.6\% \\
 & \texttt{mistralai/mistral-small-2603}   & 329 &   6 & 0 & 98.2\% \\
 & \texttt{z-ai/glm-5-turbo}               & 240 &  93 & 2 & 72.1\% \\
 & \textbf{panel-2of3}                     & \textbf{241} & \textbf{\phantom{0}92} & \textbf{2} & \textbf{72.4\%} \\
\bottomrule
\end{tabular}
\end{table}

The vote distribution underneath these aggregates further qualifies the
result. Counting accept votes among the three judges per query, the
distribution is concentrated on the split outcomes:
\texttt{qwen3.5-9b} produces $156$ $2{-}1$ accepts and $129$ $1{-}2$
rejects, with only $17$ unanimous accepts and $33$ unanimous rejects;
\texttt{deepseek-v4-flash} produces $185$ $2{-}1$ and $90$ $1{-}2$, with
$43$ unanimous accepts and $15$ unanimous rejects; \texttt{kimi2.6}
produces $193$ $2{-}1$ and $86$ $1{-}2$, with $48$ unanimous accepts and
only $6$ unanimous rejects. The monotonic shift in the unanimous-accept and
unanimous-reject counts across the three models (Qwen
$17/33$~$\to$~DeepSeek $43/15$~$\to$~Kimi $48/6$) mirrors the panel ranking
and is reassuringly robust to the choice of any single judge.

\paragraph{Reading the gap.} The single-judge canonical numbers
(Table~\ref{tab:evalresults}, $5.4$--$14.6\%$) and the panel numbers
(Table~\ref{tab:panelresults}, $51.6$--$72.4\%$) are not in tension: the
former is the strictest of the three readings (the canonical judge has the
highest threshold), the latter is the median. They are both valid summaries
of the same set of $1{,}005$ accept/reject decisions ($3$ judges $\times$
$335$ queries), differing only in how those decisions are aggregated. What
matters scientifically is that the \emph{ranking} \texttt{qwen3.5-9b} $<$
\texttt{deepseek-v4-flash} $<$ \texttt{kimi2.6} is preserved on every
judge and on the panel, so the comparative conclusion is robust even though
the absolute accuracy on this protocol is judge-dependent. We therefore
report both readings: the single-judge numbers for cross-protocol
comparability with the rest of the table, the panel numbers for a less
threshold-sensitive view of agentic planning capability.

\paragraph{Caveats.} The pairwise Cohen's $\kappa$ between
\texttt{openai/gpt-5.4-mini} and \texttt{mistralai/mistral-small-2603} on
this protocol is close to zero, reflecting the wide threshold gap and not a
random-guessing baseline -- the two judges agree on which queries are
\emph{best} (the model ranking is preserved) but disagree on where to set
the accept/reject threshold. The panel-as-median heuristic mitigates this
mechanically but does not anchor the threshold to any external ground truth;
a human-anchored $\kappa(\text{judge, human})$ over a stratified
$50$--$100$-query subset is the recommended next step.

\subsection{Aggregate Per-Query Accuracy and Oracle Bound}\label{sec:eval-oraclebound}

While Table~\ref{tab:evalresults} reports per-protocol accuracy, the routing
literature typically compares models on aggregate per-query response accuracy
and against the achievable upper bound when an oracle selects the optimal
model per query. This subsection reports those metrics on the full Dataset~A
corpus ($N{=}5{,}504$, including the $165$ \texttt{planning\_agentic\_multiturn}
queries added in \texttt{v0.4} that augment the original \texttt{v0.3.0}
release of $5{,}339$ rows).

\begin{table}[ht]
\centering
\caption{Standalone per-model response accuracy on Dataset~A
($N{=}5{,}504$). A query is \emph{solved} iff the model's response
passes the per-protocol grader. No routing is involved.}
\label{tab:standalone_response_accuracy}
\small
\begin{tabular}{l r r}
\toprule
\textbf{Model} & \textbf{Solved} & \textbf{Response accuracy} \\
\midrule
qwen3.5-9b        & 3{,}477 / 5{,}504 & 63.17\% \\
deepseek-v4-flash & 4{,}056 / 5{,}504 & 73.69\% \\
kimi2.6           & 4{,}129 / 5{,}504 & 75.02\% \\
\bottomrule
\end{tabular}
\end{table}

\begin{table}[ht]
\centering
\caption{Oracle upper bound: maximum achievable response accuracy when an
oracle selects the optimal model per query, for progressively larger model
pools.}
\label{tab:oracle_bound}
\small
\begin{tabular}{l r r}
\toprule
\textbf{Pool} & \textbf{Oracle solved} & \textbf{Oracle accuracy} \\
\midrule
kimi2.6 standalone (no routing)        & 4{,}129 / 5{,}504 & 75.02\% \\
kimi2.6 $+$ deepseek-v4-flash          & 4{,}511 / 5{,}504 & 81.96\% \\
qwen3.5-9b $+$ ds4 $+$ kimi2.6 (full)  & 4{,}582 / 5{,}504 & \textbf{83.25\%} \\
\midrule
Unsolvable by any pool model           & 922 / 5{,}504    & 16.75\% \\
\bottomrule
\end{tabular}
\end{table}

\begin{table}[ht]
\centering
\caption{Per-query solver overlap (Venn) across the three-model pool
($N{=}5{,}504$). Each row reports the (qwen3.5-9b, deepseek-v4-flash,
kimi2.6) solver status as (T, F). Bold rows highlight the routing
opportunities where a cheaper or middle-tier model rescues a kimi2.6
failure.}
\label{tab:venn_breakdown}
\small
\begin{tabular}{c c c r r l}
\toprule
\textbf{qwen} & \textbf{ds4} & \textbf{kimi} & \textbf{N} & \textbf{\%} & \textbf{Notes} \\
\midrule
T & T & T & 3{,}106 & 56.43\% & easy: all three solve \\
F & T & T & 568   & 10.32\% & ds4$+$kimi (qwen fail) \\
F & F & T & 249   & 4.52\%  & only kimi rescues \\
\textbf{F} & \textbf{T} & \textbf{F} & \textbf{288} & \textbf{5.23\%} & \textbf{only ds4 rescues (kimi fails!)} \\
T & F & T & 206   & 3.74\%  & qwen$+$kimi (ds4 fail) \\
\textbf{T} & \textbf{T} & \textbf{F} & \textbf{94}  & \textbf{1.71\%} & \textbf{qwen$+$ds4 (kimi fails!)} \\
T & F & F & 71    & 1.29\%  & only qwen rescues \\
F & F & F & 922   & 16.75\% & unsolvable (irreducible) \\
\midrule
\multicolumn{3}{r}{\textbf{Total}} & \textbf{5{,}504} & \textbf{100.00\%} & \\
\bottomrule
\end{tabular}
\end{table}

Despite \texttt{kimi2.6} leading on most protocols
(Table~\ref{tab:evalresults}), it is not the per-query oracle: it fails
$1{,}375$ queries, of which $453$ are correctly solved by at least one of the
cheaper models in the pool ($382$ by \texttt{deepseek-v4-flash} and $71$ by
\texttt{qwen3.5-9b} alone). Consequently, an oracle restricted to
(\texttt{kimi2.6}, \texttt{deepseek-v4-flash}) reaches $81.96\%$ response
accuracy -- a $6.94$~percentage-point improvement over \texttt{kimi2.6}
standalone -- and the full three-model oracle reaches $83.25\%$
(Table~\ref{tab:oracle_bound}).

The Venn breakdown (Table~\ref{tab:venn_breakdown}) bounds the achievable
routing gain. The $382$ queries where only \texttt{deepseek-v4-flash}
succeeds ($5.23\%$ of the corpus) plus the $94$ where \texttt{qwen3.5-9b}
and \texttt{deepseek-v4-flash} both succeed but \texttt{kimi2.6} fails
($1.71\%$) constitute the principal routing upside -- capturing them moves
an intelligent router strictly above any single-model baseline. The $922$
unsolvable queries ($16.75\%$) define the irreducible ceiling: closing this
gap requires extending the pool with a stronger fourth model rather than
smarter routing within the current three.

\subsection{Discussion}\label{sec:eval-truncation}

\paragraph{Reasoning-budget truncation and degenerate generation.}
An initial \texttt{qwen3.5-9b} run exhibited high \texttt{finish\_reason}~$=$~\texttt{length}
rates on \texttt{ifeval\_constraint\_check} ($63\%$), \texttt{mcq\_letter}
($27\%$), \texttt{lcb\_unit\_test} ($20\%$) and creative \texttt{rubric\_judge}
($20\%$), because the per-protocol \texttt{max\_tokens} caps were smaller than the
thinking budget plus the answer, leaving an empty (ungradable) response. Caps
were raised (Table~\ref{tab:maxtokens}) and the affected queries re-run.
Re-running resolved the empty-answer cases but exposed a second, distinct failure
mode: on $6.7\%$ of \texttt{ifeval}, $8.2\%$ of \texttt{lcb} and $9.5\%$ of
creative \texttt{rubric\_judge} queries the model enters a \emph{degenerate
generation loop}---emitting tens of thousands of characters of repeated text
(e.g.\ a code comment or a sentence repeated until the token cap). These are not
recoverable by a larger budget; they are genuine model failures and are graded as
incorrect. They account for a substantial part of \texttt{qwen3.5-9b}'s deficit
on the coding, instruction-following and creative protocols, and constitute a
robustness finding for $9$B-class reasoning models. \texttt{deepseek-v4-flash}
did not exhibit either failure mode at a material rate ($\le 1\%$ truncation
throughout). The only residual ungradable cases are 22 \texttt{math\_equiv}
queries for \texttt{qwen3.5-9b} ($4.2\%$, below the re-run threshold).

\paragraph{Kimi2.6 unbounded-thinking truncation.} \texttt{kimi2.6} exhibits a
third truncation pattern, distinct from \texttt{qwen3.5-9b}'s empty-answer and
degenerate-loop modes: a small fraction of queries cause its reasoning trace to
exceed the per-call \texttt{max\_tokens} budget on both the initial pass and a
retry pass with raised caps. The final residual ungradable rate on the
deterministic protocols is \texttt{lcb\_unit\_test} $2.4\%$,
\texttt{ifeval\_constraint\_check} $2.6\%$, and $0\%$ on \texttt{math\_equiv},
\texttt{mcq\_letter}, and \texttt{gsm8k\_final\_answer}. On the retry pass we
attempted an explicit reasoning-budget cap via OpenRouter's
\texttt{reasoning.max\_tokens} field (a $20{,}000$-token cap applied to the
reasoning-only portion): on a $12$-row pilot, the cap was honoured on only
$7/12$ queries; on the remaining $5/12$ \texttt{kimi2.6} emitted between
$25\,k$ and $41\,k$ reasoning tokens, indicating the parameter is not enforced
by this provider. We therefore fell back to raising the joint
\texttt{max\_tokens} cap alone (\texttt{lcb\_unit\_test}~$\to 65{,}536$,
\texttt{math\_equiv}~$\to 49{,}152$,
\texttt{ifeval\_constraint\_check}~$\to 32{,}768$), which recovered $142/188$
of the initial skips. The residual $46$ rows have no in-protocol mechanism to
bound the trace and are graded as None (skip-on-truncation), contributing to
the Acc$_c$ versus Acc$_o$ gap on \texttt{ifeval\_constraint\_check} and
\texttt{lcb\_unit\_test}.

\paragraph{Agentic planning is the hard frontier.} Across the three agentic
protocols, the picture is one of sharp asymmetry. Single-turn
\texttt{tool\_call\_match} is near-saturated for all three models ($93$--$94\%$),
indicating that the BFCL single-turn function-calling format is no longer
discriminative at this model scale. The two harder agentic protocols collapse
in the opposite direction: BFCL multi-turn lands in the $18$--$21\%$ band for
all three models, and planning \texttt{rubric\_judge} lands in the
$5.4$--$14.6\%$ band under the canonical single judge -- or, more
threshold-robustly, in the $51.6$--$72.4\%$ band under the three-judge panel
(Table~\ref{tab:panelresults}). Manual inspection of the planning rejections
under the canonical judge confirms substantive failure modes: required tools
left unused, explicit constraints (budget, deadline) violated, or hypothetical
and non-actionable steps substituted for concrete ones. The conjunctive rubric
(all applicable criteria must pass) compounds the difficulty. Critically, the
panel-versus-canonical gap on planning does not reorder the models: on every
judge and on the panel, \texttt{kimi2.6} $>$ \texttt{deepseek-v4-flash} $>$
\texttt{qwen3.5-9b}. The multi-turn results paint a similar weakness story
across providers ($\le 22.7\%$ conditional accuracy on a deterministic
state-based checker) and suggest that multi-step agentic tool-use, not
function-calling format, is the bottleneck at this scale. A graded
per-criterion planning score and an explicit decomposition of multi-turn
failures into tool-selection vs.\ state-tracking errors are left to future
work.

\paragraph{Short-form factuality.} On SimpleQA-style \texttt{llm\_judge\_factual},
\texttt{qwen3.5-9b} is correct on only $11.9\%$ of attempted queries, against
$49.8\%$ for \texttt{deepseek-v4-flash} and $59.6\%$ for \texttt{kimi2.6}.
Abstention behaviour, however, differs sharply across the three models:
\texttt{qwen3.5-9b} and \texttt{deepseek-v4-flash} decline only $\approx 7$--$8\%$
of queries, so their conditional/overall accuracies are close, whereas
\texttt{kimi2.6} declines $51.4\%$ of queries -- collapsing its overall accuracy
to $28.7\%$ despite the higher conditional rate. The \texttt{qwen3.5-9b} vs.\
\texttt{deepseek-v4-flash} gap is consistent with parametric-knowledge capacity
scaling with model size; the \texttt{kimi2.6} pattern is qualitatively distinct
and is analysed next.

\paragraph{Calibration and hallucination rates.}\label{sec:eval-calibration}
The three models exhibit qualitatively different calibration behaviours on the
two judge-graded protocols. Defining the \emph{hallucination rate}
$H = F / (T+F)$ on attempted queries only (i.e.\ excluding
\textsc{not\_attempted} abstentions, following the SimpleQA
convention~\cite{wei2024simpleqa}), we observe:

\begin{center}
\small
\begin{tabular}{l c c c}
\toprule
\textbf{Protocol} & \texttt{qwen3.5-9b} & \texttt{deepseek-v4-flash} & \texttt{kimi2.6} \\
\midrule
\texttt{llm\_judge\_factual}         & $88.1\%$ & $50.2\%$ & $\mathbf{40.4\%}$ \\
\texttt{rubric\_judge} (creative)    & $49.0\%$ & $34.2\%$ & $\mathbf{24.3\%}$ \\
\bottomrule
\end{tabular}
\end{center}

\noindent\texttt{kimi2.6} shows the lowest hallucination rate on both protocols,
but the pattern is most pronounced on \texttt{llm\_judge\_factual}, where it
declines $51.4\%$ of queries (\textsc{not\_attempted}) versus $7$--$8\%$ for the
other two models. Inspection of the $360$ declined responses shows
$337$ are the literal string \emph{``I don't know''}, preceded on average by
$1{,}644$ reasoning tokens (median $1{,}316$) -- i.e.\ an active refusal after
substantial deliberation, not a parse artefact. On the very same $360$ queries
\texttt{qwen3.5-9b} answers incorrectly $81.4\%$ of the time and
\texttt{deepseek-v4-flash} $53.3\%$, indicating the declined items are
genuinely difficult and that \texttt{kimi2.6}'s abstention is selective rather
than blanket. This trade-off -- lower coverage in exchange for higher
reliability when answering -- is precisely the behaviour the SimpleQA
framework~\cite{wei2024simpleqa} is designed to reward, and it suggests the
model has been calibration-trained to prefer honest abstention over fabrication.

\paragraph{Scope.} All three target models are now evaluated on all ten
protocols, including the three agentic protocols (planning
\texttt{rubric\_judge} under both the canonical single judge and the
three-judge panel, single-turn \texttt{tool\_call\_match}, and multi-turn
\texttt{tool\_call\_match}). The three-judge panel is currently applied only
to the planning rubric; extending it to the creative \texttt{rubric\_judge}
and to \texttt{llm\_judge\_factual} is methodologically straightforward and
left to follow-up. A human-anchored
$\kappa(\text{judge}, \text{human})$ on a stratified $50$--$100$-query subset
of the planning rubric is the recommended next step to validate the
panel-as-median heuristic. The routing-system evaluation is the subsequent
step.

% =============================================================================
\section{Limitations}\label{sec:limitations}
% =============================================================================

We document the following limitations.

\begin{enumerate}[leftmargin=*,itemsep=2pt]
  \item \textbf{Tokenizer proxies.} Two of the three tokenizers (DeepSeek-V3 for DeepSeek-V4-Flash; Kimi-K2.5 for Kimi-2.6) are proxies for closed-weight target models. Drift relative to the true production tokenizer is estimated at $\pm$2--5\% and is documented in the dataset card. Cost-tier estimation downstream should treat these counts as lower-bound proxies.
  \item \textbf{Gated masking.} GAIA (\cite{mialon2023gaia}) and GPQA-Diamond (\cite{rein2024gpqa}) rows are masked in the public release. Users with credentialed access can restore the original payloads via \texttt{99\_unmask\_gated.py}; routing evaluations on these subsets must be performed locally on the unmasked variant.
  \item \textbf{Custom-Validated genre clustering.} The 100 approved Custom-Validated entries exhibit a residual genre clustering toward poetic / lighthouse / ocean themes (observed in approximately 30\% of the pool). Future expansion to the full target of 404 entries should oversample under-represented genre tags.
  \item \textbf{AIME 2024 dropped.} AIME 2024 was excluded due to insufficient statistical power at $n{=}10$ post-cutoff prompts and unresolvable contamination concerns. Only AIME 2025 ($n{=}30$) is retained.
  \item \textbf{Multi-IF excluded.} The Multi-IF benchmark (cc-by-nc-2.0) was excluded because its license is incompatible with the public push. The 200 instruction-following slots vacated by Multi-IF were reallocated to additional IFEval rows (apache-2.0).
  \item \textbf{Synthetic few-shot pools.} For the six sources without a native train split, the few-shot exemplars are synthetic (generated by \texttt{qwen3.5-122b@regolo}). Routing evaluations consume this pool consistently across all three target models, so the synthetic pool affects absolute scores but not relative ranking — which is the routing-evaluation primitive of interest.
  \item \textbf{Irreducible pool ceiling.} As reported in \S\ref{sec:eval-oraclebound} and Table~\ref{tab:oracle_bound}, $922$ queries ($16.75\%$ of the corpus) are not solved by any of the three target models, which bounds the achievable response accuracy of \emph{any} router restricted to this pool at $83.25\%$. Closing the gap requires extending the pool with at least one model strictly stronger on these queries (e.g., a larger frontier model), not refining the routing logic itself.
\end{enumerate}

% =============================================================================
\section{Distribution and Reproducibility}\label{sec:distribution}
% =============================================================================

Dataset A is publicly distributed on the Hugging Face Hub under \texttt{massaindustries/dataset-A-routing-eval} (revision \texttt{73aa0a09}). The repository is laid out in the multi-config format with seven configurations: a top-level \texttt{all} configuration ($N{=}5{,}339$) plus six per-dimension configurations matching the row counts in Table~\ref{tab:composition}. The top-level license is \texttt{other}; per-source license is recorded in-row in the \texttt{license} field.

Reproducibility is supported by (i)~per-source upstream revision SHAs pinned in \texttt{data/reports/lockfile.yaml}; (ii)~SHA-256 of every synthetic few-shot pool committed alongside the generator-model identifier; (iii)~deterministic seeds in every sampling and shuffling stage; (iv)~a final sort by \texttt{query\_id} after all shuffles to remove cross-revision instability; and (v)~a smoke-test script (\texttt{80\_verify\_dataset.py}) that re-loads the published artifact, re-tokenizes a 50-row sample, and asserts schema and token-count invariants.

% =============================================================================
\begin{thebibliography}{99}
% =============================================================================

\bibitem{cohen1960kappa}
J.~Cohen.
\newblock A coefficient of agreement for nominal scales.
\newblock {\em Educational and Psychological Measurement}, 20(1):37--46, 1960.

\bibitem{hendrycks2021math}
D.~Hendrycks, C.~Burns, S.~Kadavath, A.~Arora, S.~Basart, E.~Tang, D.~Song, and J.~Steinhardt.
\newblock Measuring mathematical problem solving with the {MATH} dataset.
\newblock {\em NeurIPS Datasets and Benchmarks Track}, 2021.

\bibitem{cobbe2021gsm8k}
K.~Cobbe, V.~Kosaraju, M.~Bavarian, M.~Chen, H.~Jun, L.~Kaiser, M.~Plappert, J.~Tworek, J.~Hilton, R.~Nakano, C.~Hesse, and J.~Schulman.
\newblock Training verifiers to solve math word problems.
\newblock {\em arXiv preprint arXiv:2110.14168}, 2021.

\bibitem{wang2024mmlupro}
Y.~Wang, X.~Ma, G.~Zhang, Y.~Ni, A.~Chandra, S.~Guo, W.~Ren, A.~Arulraj, X.~He, Z.~Jiang, T.~Li, M.~Ku, K.~Wang, A.~Zhuang, R.~Fan, X.~Yue, and W.~Chen.
\newblock {MMLU}-{Pro}: A more robust and challenging multi-task language understanding benchmark.
\newblock {\em arXiv preprint arXiv:2406.01574}, 2024.

\bibitem{zhou2023ifeval}
J.~Zhou, T.~Lu, S.~Mishra, S.~Brahma, S.~Basu, Y.~Luan, D.~Zhou, and L.~Hou.
\newblock Instruction-following evaluation for large language models.
\newblock {\em arXiv preprint arXiv:2311.07911}, 2023.

\bibitem{pyatkin2025ifbench}
V.~Pyatkin et~al.
\newblock {IFBench}: An instruction-following benchmark for {LLM} robustness evaluation.
\newblock {\em AllenAI Technical Report}, 2025.

\bibitem{yan2024bfcl}
F.~Yan, H.~Mao, C.~C.-J. Ji, T.~Zhang, S.~G. Patil, I.~Stoica, and J.~E. Gonzalez.
\newblock Berkeley function calling leaderboard.
\newblock \url{https://gorilla.cs.berkeley.edu/leaderboard.html}, 2024.

\bibitem{jain2024livecodebench}
N.~Jain, K.~Han, A.~Gu, W.-D. Li, F.~Yan, T.~Zhang, S.~Wang, A.~Solar-Lezama, K.~Sen, and I.~Stoica.
\newblock {LiveCodeBench}: Holistic and contamination free evaluation of large language models for code.
\newblock {\em arXiv preprint arXiv:2403.07974}, 2024.

\bibitem{wei2024simpleqa}
J.~Wei, N.~Karina, H.~W. Chung, Y.~J. Jiao, S.~Papay, A.~Glaese, J.~Schulman, and W.~Fedus.
\newblock {SimpleQA}: Measuring short-form factuality in large language models.
\newblock {\em OpenAI Technical Report}, 2024.

\bibitem{rein2024gpqa}
D.~Rein, B.~L. Hou, A.~C. Stickland, J.~Petty, R.~Y. Pang, J.~Dirani, J.~Michael, and S.~R. Bowman.
\newblock {GPQA}: A graduate-level google-proof {Q}\&{A} benchmark.
\newblock {\em COLM}, 2024.

\bibitem{mialon2023gaia}
G.~Mialon, C.~Fourrier, C.~Swift, T.~Wolf, Y.~LeCun, and T.~Scialom.
\newblock {GAIA}: A benchmark for general {AI} assistants.
\newblock {\em arXiv preprint arXiv:2311.12983}, 2023.

\bibitem{paech2024eqbench}
S.~J. Paech.
\newblock {EQ}-{Bench}: An emotional intelligence benchmark for large language models.
\newblock {\em arXiv preprint arXiv:2312.06281}, 2024.

\bibitem{stanford2025litbench}
Stanford CRFM.
\newblock {LitBench}: A benchmark for long-form literary writing evaluation.
\newblock \url{https://huggingface.co/datasets/SAA-Lab/LitBench-Test}, 2025.

\bibitem{yao2024taubench}
S.~Yao, N.~Shinn, P.~Razavi, and K.~Narasimhan.
\newblock $\tau$-bench: A benchmark for tool-agent-user interaction in real-world domains.
\newblock {\em arXiv preprint arXiv:2406.12045}, 2024.

\bibitem{chen2023frugalgpt}
L.~Chen, M.~Zaharia, and J.~Zou.
\newblock {FrugalGPT}: How to use large language models while reducing cost and improving performance.
\newblock {\em arXiv preprint arXiv:2305.05176}, 2023.

\bibitem{ong2024routellm}
I.~Ong, A.~Almahairi, V.~Wu, W.-L. Chiang, T.~Wu, J.~E. Gonzalez, M.~W. Kadous, and I.~Stoica.
\newblock {RouteLLM}: Learning to route {LLM}s with preference data.
\newblock {\em arXiv preprint arXiv:2406.18665}, 2024.

\end{thebibliography}

\end{document}
